%! Factoring Trinomials — Unit 1, Lesson 1.4 — Warm-Up (Answer Key)
\documentclass[10pt]{article}
\usepackage{atda-article}
\usepackage{atda-key}   % pulls in -boxes; do NOT also load -boxes
\newcommand{\TallMath}[1]{$\displaystyle #1\rule[-1.4em]{0pt}{3.2em}$}
\setlength{\workrowsep}{6pt}   % handwriting room; moves blank and key together

\begin{document}

\pageheader{Unit 1, Lesson 1.4}{Warm-Up --- Answer Key}
% NAMESTRIP: no \namedateperiod here — mirrors the blank. Teacher-only prose goes
% in the lesson plan as \begin{teachernote}[Warm-Up], never in a key.

\vspace{0.15in}

\begin{notesbox}{Four Moves Today's Lesson Runs On}
\small
Item 4 uses your answers from items 2 and 3 --- do them in order. Show your work on all four.

\begin{enumerate}[leftmargin=1.6em, itemsep=6pt, topsep=4pt]

\item \textbf{Spiral --- Lesson 1.1.} Multiply out: \quad (a) $(x+4)(x+7)$ \quad
(b) $(2x+3)(x-5)$
\begin{work}
  (x+4)(x+7) &= x^2+11x+28 \\
  (2x+3)(x-5) &= 2x^2-7x-15
\end{work}

\item \textbf{Spiral --- Lesson 1.3.} Factor by grouping: $2x^2+6x+5x+15$
\begin{work}
  2x^2+6x+5x+15 &= 2x(x+3) + 5(x+3) \\
                &= (x+3)(2x+5)
\end{work}

\item \textbf{Number pairs.} Find two integers with the given product and sum:
\quad product $12$, sum $7$; \quad product $-18$, sum $3$; \quad product $24$, sum $-11$
\begin{work}
  12, \; 7 &\Rightarrow 3 \text{ and } 4 \qquad -18, \; 3 \Rightarrow 6 \text{ and } -3
           \qquad 24, \; -11 \Rightarrow -3 \text{ and } -8
\end{work}

\item \textbf{Discovery.} The four terms in item 2 came from the trinomial $2x^2+11x+15$ --- the
middle term $11x$ was split into $6x+5x$. Compare the split numbers $6$ and $5$ to the
trinomial's $2$ and $15$. What do you notice about their sum, and about their product?
\begin{work}
  6+5 &= 11 \text{, the middle coefficient} \\
  6 \cdot 5 &= 30 = 2 \cdot 15 \text{, the leading coefficient times the constant}
\end{work}

\end{enumerate}
\end{notesbox}

\end{document}
